% Baccalauréat malgache, série D, 2004 : énoncé.
\TCDBacAttribution{D}{2004}
\providecommand{\Courbe}{(\mathscr{C})}
\providecommand{\D}{(\mathscr{D})}
\providecommand{\Plan}{(\mathscr{P})}
\providecommand{\Ouv}{(O ; \vec{u} , \vec{v})}
%\providecommand{\Oij}{(O ; \vec{\imath} , \vec{\jmath})}
\providecommand{\cm}{~\mathrm{cm}}
\providecommand{\dx}{~\mathrm{d}x}
\providecommand{\Arg}[1]{\mathrm{Arg}\left( #1 \right)}
\providecommand{\Angl}[2]{\left( \overrightarrow{#1} ; \overrightarrow{#2} \right)}
%%%%%%%%%%%%%%%%%%%%%%%%%%%%%%%%%%%%%%%%%%%%%%%%%%%%%%%%%%%%%%%%%%%%%%%%%%%%%%%%%%%%%%%%%%%%%%%
%%%%%%%%%%%%%%%%%%%                   %%%%%%%%%%%%%%%%%%%                   %%%%%%%%%%%%%%%%%%%
%%%%%%%%%%%%%%%%%%%%%%%%%%%%%%%%%%%%%%%%%%%%%%%%%%%%%%%%%%%%%%%%%%%%%%%%%%%%%%%%%%%%%%%%%%%%%%%
\noindent \underline{\textsc{Exercice I}} (5 points)
      \begin{enumerate}
            \item R\'esoudre dans $\mathbb{C}$ l'\'equation d'inconnue $z$ : $z^2 - (4 + 5i) \; z - 1 + 7 i = 0$.
            \item Le plan complexe $\Plan$ est muni d'un rep\`ere orthonorm\'e $\Ouv$ (unit\'e $1 \cm$).
                  On consid\`ere les points $A$, $B$ et $C$ d'affixes respectives : 
                  $z_A = 1 + i$ ; $z_B = 3 + 4 i$ et $z_C = 4 - i$.
                  \begin{enumerate}
                        \item D\'eterminer l'expression complexe associ\'ee \`a la similitude plane directe $S$ telle que 
                              $S(A) = B$ et $S(B) = C$.
                        \item Pr\'eciser les \'el\'ements caract\'eristiques de $S$.
                  \end{enumerate}
            \item On note par $I$ le milieu du segment $[BC]$.
                  \begin{enumerate}
                        \item Calculer l'affixe $z_I$ de $I$.
                        \item Placer les points $A$, $B$, $C$ et $I$ dans le plan complexe $\Plan$.
                        \item D\'eterminer et construire l'ensemble $\Courbe$ des points $M$ d'affixe $z$ v\'erifiant :
                              $\displaystyle{
                                    \left|
                                          z - \frac{7 + 3 i}{2}
                                    \right|
                                           = \frac{\sqrt{26}}{2} 
                              }$.
                              D\'eterminer et construire l'ensemble $\CourbeP$ image de $\Courbe$ par $S$.
                  \end{enumerate}
      \end{enumerate}
%%%%%%%%%%%%%%%%%%%%%%%%%%%%%%%%%%%%%%%%%%%%%%%%%%%%%%%%%%%%%%%%%%%%%%%%%%%%%%%%%%%%%%%%%%%%%%%
%%%%%%%%%%%%%%%%%%%                   %%%%%%%%%%%%%%%%%%%                   %%%%%%%%%%%%%%%%%%%
%%%%%%%%%%%%%%%%%%%%%%%%%%%%%%%%%%%%%%%%%%%%%%%%%%%%%%%%%%%%%%%%%%%%%%%%%%%%%%%%%%%%%%%%%%%%%%%
\noindent \underline{\textsc{Exercice II}} (5 points) \newline
      Une urne $U_1$ contient six boules : 
      une boule num\'erot\'ee 0, deux boules num\'erot\'ees 1 et trois boules num\'erot\'ees 2.
      Une autre urne $U_2$ contient cinq boules : 
      deux boules num\'erot\'ees 0, une boule num\'erot\'ee 1 et deux boules num\'erot\'ees 2.
      Les boules sont indiscernables au toucher.
      \begin{enumerate}
            \item On tire au hasard et simultan\'ement trois boules de l'urne $U_1$.
                  Calculer les probabilit\'es des \'ev\`enements suivants : 
                  \begin{itemize}
                        \item $A$ : << La somme des num\'eros not\'es est \'egale \`a 4 >>.
                        \item $B$ : << Parmi les trois boules tir\'ees, deux boules exactement sont num\'erot\'ees par 2 >>.
                        \item $C$ : << Le produit des num\'eros not\'es sur les trois boules tir\'ees est diff\'erent de 0 >>.
                        \item $D$ : << Les trois boules obtenues portent le m\^eme num\'ero >>.
                  \end{itemize}
            \item On remet l'urne $U_1$ dans son \`etat initial.
                  On tire une boule de l'urne $U_1$, puis on tire au hasard et simultan\'ement deux boules de l'urne $U_2$.
                  On suppose que les \'ev\`enements \'el\'ementaires sont \'equiprobables.
                  Soit $X$ la variable al\'eatoire qui, \`a chaque \'eventualit\'e, associe le produit des num\'eros
                  not\'es sur les trois boules obtenues.
                  \begin{enumerate}
                        \item V\'erifier que l'univers image de $X$ est \'egal \`a l'ensemble $\{ 0, 2, 4, 8 \}$.
                        \item Montrer que $P(X = 0) = \frac{3}{4}$ et $P(X \geq 4) = \frac{11}{60}$.
                        \item Calculer la probabilit\'e de l'\'ev\`enement : $(X = 2)$.
                        \item Compl\'eter le tableau des valeurs ci-dessous :
                              \[
                                    \begin{array}{|l|r|r|r|r|}
                                          \hline
                                                k         & 0           & 2       & 4      & 8      \\
                                          \hline
                                                P(X = k)  & \frac{3}{4} &         &        &        \\
                                          \hline
                                    \end{array}
                              \]
                        \item Montrer que $E(X) = \frac{16}{15}$, $E(X)$ \'etant l'esp\'erance math\'ematique de 
                              la variable al\'eatoire $X$.
                        \item D\'efinir et repr\'esenter graphiquement la fonction de r\'epartition $F$ de $X$.
                              On utilisera pour la repr\'esentation graphique de $F$ un rep\`ere orthogonal \Oij .
                              (On prendra comme unit\'e : $1 \cm$ sur l'axe des abscisses et $6 \cm$ sur l'axe des ordonn\'ees.)
                              
                              (\textbf{N.B :} On donnera les r\'esultats sous forme de fraction irr\'eductible.)
                  \end{enumerate}
      \end{enumerate}
%%%%%%%%%%%%%%%%%%%%%%%%%%%%%%%%%%%%%%%%%%%%%%%%%%%%%%%%%%%%%%%%%%%%%%%%%%%%%%%%%%%%%%%%%%%%%%%
%%%%%%%%%%%%%%%%%%%                   %%%%%%%%%%%%%%%%%%%                   %%%%%%%%%%%%%%%%%%%
%%%%%%%%%%%%%%%%%%%%%%%%%%%%%%%%%%%%%%%%%%%%%%%%%%%%%%%%%%%%%%%%%%%%%%%%%%%%%%%%%%%%%%%%%%%%%%%
\noindent \underline{\textsc{Probleme}} (10 points) \newline
      Soit la fonction num\'erique $f$ d\'efinie sur l'intervalle $]0;+\infty[$ par : 
      $\displaystyle{f(x) = \expo{x} \ln x + \frac{\expo{x}}{x}}$.
      On note par $\Courbe$ sa courbe repr\'esentative dans un rep\`ere orthonorm\'e \Oij{} (unit\'e : $2 \cm$).
      \begin{enumerate}
            \item On donne la fonction num\'erique $g$ d\'efinie sur l'intervalle $]0;+\infty[$ par : 
                  $\displaystyle{g(x) = \ln x + \frac{2}{x} - \frac{1}{x^2}}$.
                  \label{quest: Bac D 2004 Pbm quest 1}
                  \begin{enumerate}
                        \item Calculer $\displaystyle{\lim_{x\rightarrow 0^+} g(x)}$ et 
                              $\displaystyle{\lim_{x\rightarrow +\infty} g(x)}$.
                        \item Pour tout $x > 0$, calculer $g'(x)$ et \'etudier son signe.
                        \item Dresser le tableau de variation de $g$.
                        \item Montrer qu'il existe un et un seul nombre r\'eel $\alpha$ v\'erifiant $g(\alpha) = 0$
                              avec $\frac{1}{2} < \alpha < 1$.
                        \label{quest: Bac D 2004 Pbm quest 1/ d)}
                  \end{enumerate}
            \item \begin{enumerate}
                        \item Pour tout $x > 0$, calculer $f'(x)$.
                        \item Pour tout $x > 0$, montrer l'\'egalit\'e $f'(x) = \expo{x} \; g(x)$.
                        \item Calculer $f'(\alpha)$.
                        \item Calculer $\displaystyle{\lim_{x\rightarrow 0^+} f(x)}$ et 
                              $\displaystyle{\lim_{x\rightarrow +\infty} f(x)}$.
                              \Big( 
                                    On pourra utiliser l'\'egalit\'e $\displaystyle{f(x) = \frac{\expo{x}}{x} (x \ln x + 1)}$.
                              \Big)
                        \item Dresser le tableau de variation de la fonction $f$.
                  \end{enumerate}
            \item \begin{enumerate}
                        \item Donner une \'equation de la tangente $(T)$ \`a $\Courbe$ au point d'abscisse $1$.
                        \item Construire $\Courbe$ et $(T)$ dans un m\^eme rep\`ere.
                  \end{enumerate}
            \item On pose pour tout $x > 0$, $h(x) = \expo{x} \ln x$.
                  Montrer que $h$ est une primitive de $f$ sur l'intervalle $]0 ; +\infty[$.
            \item Soit $I = \int_{\alpha}^{1} f(x) \dx$ o\`u $\alpha$ est le r\'eel d\'efini \`a la question 
                  \ref{quest: Bac D 2004 Pbm quest 1}~\ref{quest: Bac D 2004 Pbm quest 1/ d)}.
                  \begin{enumerate}
                        \item Montrer que 
                              $\displaystyle{ I = \expo{\alpha} \left( \frac{2}{\alpha} - \frac{1}{\alpha^2} \right) }$.
                        \item Interpr\'eter g\'eom\'etriquement l'int\'egrale $I$.
                  \end{enumerate}
                  On donne $\ln 2 \simeq 0,7$, $\alpha \simeq 0,6$, $f(\alpha) \simeq 2,1$, $e \simeq 2,7$.
      \end{enumerate}






